\begin{exbox}[窄带高斯波包与三阶相位余项]
\label{sec:feqo-dispersion-expansion}
取中心动量 $\mathbf p_0=p_0\mathbf e_z$，正能色散的完整二阶展开为
\begin{equation}
E(\mathbf p)=E_0+v_0\delta p_z+
\frac{\delta p_x^2+\delta p_y^2}{2\gamma_0m_{\rm e}}
+\frac{\delta p_z^2}{2\gamma_0^3m_{\rm e}}+R_3(\delta\mathbf p).
\label{eq:feqo-dispersion-hessian}
\end{equation}
这由 $\partial_i\partial_jE=c^2\delta_{ij}/E-c^4p_ip_j/E^3$ 在中心处取值。
固定正能自旋，把纵向动量包络取为
\begin{equation}
a(p)=(2\pi\sigma_p^2)^{-1/4}
\exp[-(p-p_0)^2/(4\sigma_p^2)].
\label{eq:feqo-gaussian-momentum}
\end{equation}
在窄带 $\sigma_p\ll p_0$ 内，负向尾概率
$\tfrac12\operatorname{erfc}(p_0/(\sqrt2\sigma_p))$ 可忽略。
令 $m_\parallel=\gamma_0^3m_{\rm e}$，
使用前面色散 Hessian 的纵向分量，并令 $\xi=z-v_0t$。
抽去共同载波后，Fourier 积分中的指数为
\[
-\frac{\delta p^2}{4\sigma_p^2}
\left(1+\ii\frac{2\sigma_p^2t}{m_\parallel\hbar}\right)
+\frac{\ii\delta p\xi}{\hbar}.
\]
定义 $\sigma_{z0}=\hbar/(2\sigma_p)$、
$\tau=\hbar t/(2m_\parallel\sigma_{z0}^2)$。
配方并完成高斯积分，得到归一化包络
\begin{align}
\psi(\xi,t)&=(2\pi\sigma_{z0}^2)^{-1/4}(1+\ii\tau)^{-1/2}
\exp\left[-\frac{\xi^2}{4\sigma_{z0}^2(1+\ii\tau)}\right],\label{eq:feqo-gaussian-prop-result}\\
|\psi(\xi,t)|^2&=
\frac{\exp[-\xi^2/(2\sigma_{z0}^2(1+\tau^2))]}
{\sqrt{2\pi}\sigma_{z0}\sqrt{1+\tau^2}},\label{eq:feqo-gaussian-density}\\
\sigma_z(t)^2&=\sigma_{z0}^2+
\frac{\hbar^2t^2}{4m_\parallel^2\sigma_{z0}^2}.\label{eq:feqo-gaussian-spreading}
\end{align}
$\tau=1$ 时方差增至两倍，标准差增至 $\sqrt2$ 倍。
对 $\SI{200}{\kilo\electronvolt}$ 电子，$\sigma_{z0}=\SI{1}{\nano\metre}$
给此时间 $t_*=\SI{4.65e-14}{\second}$、漂移长度 $v_0t_*=\SI{9.70}{\micro\metre}$；
初宽增至 $\SI{10}{\nano\metre}$ 时，这两个尺度都增大一百倍。

高斯没有有限支集。实际误差必须同时计入主体内的相位余项和主体外的尾部。
取 $|\delta p|\le r\sigma_p$，令 $M_3$ 是此窗内 $|E'''|$ 的上界。
尾概率为 $\eta_r=\operatorname{erfc}(r/\sqrt2)$，则
\begin{equation}
\|(U_{\rm ex}-U_2)\psi_0\|
\le\frac{tM_3(r\sigma_p)^3}{6\hbar}+2\sqrt{\eta_r}.
\label{eq:feqo-gaussian-propagation-tail-control}
\end{equation}
主体界由 $|\ee^{\ii x}-1|\le|x|$，尾界由两个幺正算符的范数和给出。
若旋量随动量变化，还需独立加入冻结旋量的幅度误差
$\sup_{|\delta p|\le r\sigma_p}\|u_s(p)-u_s(p_0)\|$。
因此窄带模型的适用范围由累计谱相位、尾概率及旋量变化共同确定。
\end{exbox}

\begin{figure}[htbp]
\centering
\includegraphics[width=0.76\linewidth]{relativistic_effective_mass.pdf}
\caption{同一相对论色散的横向与纵向曲率质量比，分别为 $\gamma$ 和 $\gamma^3$。}
\label{fig:relativistic-effective-mass}
\end{figure}
\begin{figure}[htbp]
\centering
\includegraphics[width=0.76\linewidth]{dispersion_control.pdf}
\caption{200\,\si{\kilo\electronvolt} 电子的精确色散与局部一阶、二阶展开。高阶残差还需乘传播时间除以 $\hbar$，才能判断传播相位误差。}
\label{fig:dispersion-control-modern}
\end{figure}
